\bookchapter{제22장: AI는 스마트폰에서 어떻게 작동하는가}

\chapterintro{스마트폰의 가장 멋진 “마술 현장”에 오신 것을 환영합니다, 주인공은: 인공지능. 매일 사용하시는 “찍기만 하면 나오는” 사진, 입을 열면 바로 응답하는 음성 비서, 사용할수록 더 잘 이해하는 입력기, 모두 이 보이지 않는 두뇌가 뒤에 있습니다.하나, AI 카메라: “후기 작업실”이 스마트폰으로 들어왔다현재의 AI 계산 사진술은 전문 보정사 + 사진작가 팀을 셔터 버튼에 직접 넣은 것과 비슷합니다. 셔터를 누르기 전에 AI는 이미 작업을 시작합니다: 장면을 스캔하고 “이건 파란 하늘”, “이건 얼굴”, “이건 야경”이라고 식별합니다.둘, 음성 비서: “음성 비서”는 어떻게 당신을 이해하는가?\textbf{명확하게 듣기(음성 인식)}: 음파를 텍스트로 변환.\textbf{이해하기(자연어 이해)}: 의도를 이해하고 핵심 개체를 추출.\textbf{실행 및 응답}: 해당 앱이나 서비스를 호출한 후, 음성 합성으로 “말”하여 들려줍니다.셋, 온디바이스 AI vs. 클라우드 AI온디바이스 AI는 계산을 스마트폰 안에 유지하여, 빠르고, 프라이버시가 좋으며, 오프라인에서도 작동합니다. 클라우드 AI는 인터넷 연결이 필요하지만, 더 복잡한 작업을 처리할 수 있습니다.넷, AI 대형 모델과 스마트폰: 당신의 휴대용 스마트 어시스턴트2025-2026년, AI 대형 모델(예: DeepSeek R1, ChatGPT-5, 원신이옌 5.0, 통이치엔원 3 등)의 폭발적인 성장이 스마트폰 사용 방식을 깊이 변화시켰습니다. 스마트폰은 더 이상 명령을 실행하는 도구가 아니라, 이해하고, 추론하고, 창작하는 스마트 파트너입니다. 온디바이스 AI 칩(NPU 연산 능력이 80 TOPS 돌파)이 대형 모델 추론을 처음으로 스마트폰 로컬에서 실행할 수 있게 하여, 인터넷 없이도 AI를 사용할 수 있습니다.AI 대형 모델이란 무엇인가?AI 대형 모델(Large Language Model, LLM)은 방대한 텍스트 데이터로 학습된 AI 시스템으로, 자연어를 이해하고 생성할 수 있습니다. 인터넷의 거의 모든 텍스트를 읽은 “슈퍼 학업 천재”라고 상상할 수 있으며, 질문에 답하고, 글을 쓰고, 번역하고, 프로그래밍하고, 데이터를 분석할 수 있습니다.스마트폰의 AI 대형 모델 응용:\textbf{대화 어시스턴트}: DeepSeek, ChatGPT, 더우바오, Kimi 등의 앱으로, 언제든지 질문하고, 번역하고, 카피를 쓰고, 기사를 요약할 수 있습니다. 전통적인 검색 엔진보다 더 스마트합니다 — 링크만 제공하는 것이 아니라 직접 답변을 제공합니다.\textbf{글쓰기 및 창작}: AI가 이메일 작성, 회의록 요약, 글 다듬기, 크리에이티브 카피 생성을 도와줍니다.\textbf{이미지 생성 및 편집}: AI가 한 번에 사진의 보행자 제거, 이미지 경계 확장, 스타일 변환, 심지어 텍스트 설명에 따른 이미지 생성.\textbf{실시간 번역}: 통화 실시간 번역, 대면 대화 번역, 사진 번역으로, 해외 여행에서 언어 장벽이 없습니다.\textbf{학습 과외}: AI가 개념을 설명하고, 연습 문제를 내고, 숙제를 채점할 수 있어, 휴대용 과외입니다.온디바이스 대형 모델: 스마트폰 자체로 AI 실행최신 플래그십 스마트폰은 이미 로컬에서 소형 AI 대형 모델(보통 7B 파라미터 정도)을 실행할 수 있으며, 가져온 변화는:\textbf{네트워크 불필요}: 인터넷 없이도 AI 기능 사용 가능, 비행기, 지하철에서도 가능.\textbf{프라이버시 더 안전}: 대화와 데이터가 클라우드로 업로드되지 않으며, 스마트폰 로컬에서 완전히 처리됩니다.\textbf{응답 더 빠름}: 네트워크 지연 없이, AI 응답이 거의 즉시.\textbf{현재 제한}: 온디바이스 모델의 능력은 여전히 클라우드 대형 모델에 미치지 못하며, 복잡한 작업은 여전히 인터넷이 필요합니다. 하지만 칩 AI 연산 능력의 향상으로, 온디바이스 모델은 점점 더 강력해질 것입니다.시스템 수준 AI 통합:\textbf{Apple Apple Intelligence}: iOS 18에서 도입된 AI 시스템으로, 글쓰기 도구, 이미지 매직 완드, Siri 강화, 알림 요약 기능을 통합. 온디바이스 처리 우선, 복잡한 작업만 클라우드 사용.\textbf{화웨이 Pangu 대형 모델}: HarmonyOS 시스템에 통합되어, XiaoYi 어시스턴트가 대형 모델 능력을 얻었으며, 문서 요약, 스마트 화면 인식, AI 사진 편집 등을 지원.\textbf{샤오미 Surge AI}: 시스템 수준 AI 능력으로, AI 사진 검색, AI 요약, AI 실시간 자막, 온디바이스 대형 모델 대화 등을 지원.\textbf{삼성 Galaxy AI}: 통화 실시간 번역, 노트 AI 요약, 이미지 생성 및 편집, 검색 동그라미 선택 등 기능.다섯, 몇 가지 일반적인 응용 시나리오와 실습스마트폰 AI의 “초능력”을 더 직관적으로 느낄 수 있도록, 아래에서 다섯 가지 구체적인 시나리오를 통해 AI가 일상생활에서 어떻게 “도우미”가 되는지 보여드립니다:시나리오 1: 실시간 번역 — 해외 여행의 휴대용 통역사\textbf{실제 성능}: 파리에서 주문할 때 프랑스어를 할 줄 모른다고요? iPhone의 “번역” 앱은 대면 대화 모드를 지원하여, 양쪽이 각자의 언어 자막을 보며, 동시 통역도 가능합니다. 안드로이드 진영의 삼성 Galaxy AI 통화 실시간 번역은 전화에서 직접 중국어로 말하면, 상대방은 영어를 듣고, 그 반대도 마찬가지입니다.\textbf{배경 원리}: 음성 인식(ASR) 기계 번역(NMT) 음성 합성(TTS) 3단계 파이프라인, 종단 간 지연이 1초 이내로 제어 가능.시나리오 2: AI 제거 — 행인에게 “망가진” 관광 사진 구하기\textbf{실제 성능}: 명소에서 멋진 셀카를 찍었는데, 배경에 다른 관광객이 가득 차 있다고요? “마법 제거”로 행인을 동그라미 표시하면, 스마트폰이 자동으로 행인을 식별하여 “지워버리고”, 배경을 지능적으로 채우며, 효과가 종종 놀랍습니다.\textbf{배경 원리}: 이미지 분할(Inpainting) 모델이 먼저 선택 영역을 식별하고 삭제한 다음, 주변 픽셀과 의미 이해에 기반하여 합리적인 채우기 내용을 “상상”합니다. 그 배경에 있는 AI 모델 학습 데이터량은 수백만 장의 이미지에 달합니다.시나리오 3: AI 사진 편집 — 한 번에 “사진 마스터”로\textbf{실제 성능}: 흐린 날의 회색 풍경 사진을 찍었나요? AI “마법 하늘”이 한 번에 하늘을 맑은 하늘, 저녁 노을, 심지어 별 하늘로 바꿀 수 있습니다. AI 이미지 확장은 세로 구도를 가로 구도로 바꿀 수 있으며, 스마트폰이 화면 밖의 장면을 “상상”합니다.\textbf{배경 원리}: 생성형 AI(Generative AI)의 확산 모델(Diffusion Model)에 기반하여, 이미지의 의미 내용을 이해한 후 새로운 픽셀을 생성합니다. 이것이 현재 AI 이미지 처리의 “최전선”입니다.시나리오 4: 문서 처리 — 사무직의 효율 배가\textbf{실제 성능}: 회의가 끝난 후 회의 녹음을 AI에 던지면, 자동으로 구조화된 회의록을 정리하고, 할 일을 추출하며, 요약을 생성합니다. 10페이지짜리 영문 PDF 보고서도 AI가 10초 내에 핵심 결론을 추출해 줍니다.\textbf{배경 원리}: 장문 맥락 대형 모델(Long Context LLM)이 수만 자의 문서를 “읽고”, 다단계 추론과 정보 추출을 수행할 수 있습니다. Apple Intelligence, 화웨이 XiaoYi, WPS AI, 텐센트 Yuanbao 등이 모두 이 능력을 갖추고 있습니다.시나리오 5: AI 사진 검색 및 물체 인식 — 당신의 “시각 검색 엔진”\textbf{실제 성능}: 거리에서 모르는 꽃을 보고, 카메라를 켜면 AI가 즉시 “이것은 무궁화입니다”라고 알려줍니다. 앨범에서 작년 바다에서 찍은 사진을 찾을 때, “바다”로 검색하면 AI가 자동으로 바다를 포함한 모든 사진을 찾아주며, 수동으로 뒤질 필요가 없습니다.\textbf{배경 원리}: 컴퓨터 비전(CV) 모델이 이미지 내용을 벡터화(Embedding)한 다음, 텍스트 쿼리와 유사도 매칭을 수행합니다. 이러한 “다중 모달 검색”은 스마트폰 앨범 관리의 큰 도약입니다.여섯, 프롬프트 최적화: AI를 당신의 지휘에 따르게 하는 “소통 공식”AI의 답변 품질은 큰 부분에서 “질문하는 방법”에 달려 있습니다. 같은 질문이라도 물어보는 방식에 따라 AI의 답변이 하늘과 땅 차이입니다. 아래에서\textbf{가장 실용적이고 가장 쉽게 시작할 수 있는 “역할-배경-작업-제한” 소통 사고방식}을 소개하고, 많은 예시를 함께 제공합니다.왜 프레임워크가 필요한가?많은 사람들이 AI를 Baidu처럼 사용합니다 — 몇 개의 키워드만 입력합니다. 하지만 AI는 검색 엔진이 아니라, “독해 후 답변”하는 언어 모델입니다.\textbf{더 많은 맥락을 제공하고, 원하는 것을 명확히 하며, 빗나가지 않도록 제한하면}, 정확한 결과를 얻을 수 있습니다. 프레임워크는 머릿속의 “좋은 답변의 모습”을 AI가 이해할 수 있는 지시로 번역하는 것입니다.1. 역할(Role): AI에게 자신이 누구인지 알려주기AI에 구체적인 “정체성”을 부여하면, 해당 분야의 전문 지식과 표현 방식을 호출합니다.\textbf{역할 설정이 구체적일수록, 출력도 더 전문적입니다}.실용 예시:“당신은\textbf{10년 경력의 샤오미 스마트폰 수리 엔지니어}입니다.”(기술 문제)“당신은\textbf{복잡한 개념을 쉽게 설명하는 중학교 물리 선생님}입니다.”(과학 문제)“당신은\textbf{심미안이 좋고, 많이 읽고, 자신만의 스타일이 있는 도서 편집자}입니다.”(글쓰기 문제)“당신은\textbf{매일 직장인에게 아침을 만드는 요리 블로거}입니다.”(생활 문제)“당신은\textbf{인내심 있고, 온화하며, 초보자를 절대 비웃지 않는 중의사}입니다.”(건강 문제)\textbf{왜 효과적인가?}역할은 AI의\textbf{지식 범위},\textbf{표현 스타일},\textbf{전문적 깊이}를 설정합니다. 역할 없는 AI는 “다 알지만 뭐 하나 잘하는 것이 없는” 일반인이고, 역할이 있으면 해당 전문가처럼 사고할 수 있습니다.2. 배경(Background): AI에게 상황을 알려주기“당신이 누구인지, 현재 상황이 어떤지, 제약 조건이 무엇인지”를 미리 명확히 말하면, AI는\textbf{정말 당신에게 적합한}답변을 줄 수 있으며, 두루뭉술한 답변을 주지 않습니다.실용 예시:“내 스마트폰은\textbf{샤오미 14, 반 년 전에 구매, 최근 게임할 때 심하게 뜨거워집니다}.”(기기 상황)“저는\textbf{60세 은퇴한 아주머니, 눈이 좋지 않고, 스마트폰을 잘 사용하지 못합니다}.”(사용자 프로필)“저는\textbf{2026년 6월 영어 6급 시험 준비 중, 기초가 보통}입니다.”(학습 배경)“저는\textbf{3선 도시에 작은 카페를 운영, 월 매출 3만, 직원 2명 고용}.”(비즈니스 배경)“저는\textbf{아기를 낳고 산후조리 중, 수면이 심각하게 부족}합니다.”(생활 단계)\textbf{왜 효과적인가?}배경 정보는 AI가\textbf{“누구를 위해 서비스하는지”}와\textbf{“어떤 조건에서”}를 알 수 있게 합니다. 같은 문제라도 학생에게 설명하는 것과 부모에게 설명하는 것은 다르고, 도시인과 농촌인에게 조언하는 것도 다릅니다. 배경이 없으면, AI는 “최대 공약수”의 답변만 줄 수 있으며, 종종 진짜 고충을 짚지 못합니다.3. 작업(Task): AI에게 무엇을 할지 명확히 알려주기이 단계에서는\textbf{동사로 시작}(쓰다, 나열하다, 대비하다, 번역하다, 설명하다, 요약하다 등)하여, 작업을\textbf{구체적이고 실행 가능하게}설명해야 합니다. 모호한 작업은 모호한 답변만 얻습니다.실용 예시:“\textbf{내가 즉시 닫아야 할 권한 5개를 나열하고}, 이유를 설명해주세요.”(동작 명확)“\textbf{다음 영어 이메일을 중국어로 번역}하되, 비즈니스 어조를 유지해주세요.”(작업 구체적)“\textbf{이 4대의 3000원대 스마트폰을 대비}하여, 장단점 대비 표를 제공해주세요.”(형식 명확)“\textbf{초등학생도 이해할 수 있는 말로 '5G가 무엇인지'를 설명}해주세요.”(스타일 명확)“\textbf{내 아침 식당에 세 가지 짧은 영상 카피를 써주세요}, 직장인, 커플, 아이 있는 부모 각각 대상으로.”(세분화 명확)\textbf{왜 효과적인가?}작업 부분은 AI 행동의 “지령 중추”입니다.\textbf{동사가 명확할수록, 요구가 구체적일수록, AI는 “무엇을 전달할지” 알게 됩니다}. 모호한 “만들어 봐”, “처리해”, “살펴봐”는 AI의 큰 금기로 — AI가 의도를 추측하다가 종종 빗나갑니다.4. 제한(Constraints): AI에게 하지 말아야 할 것 알려주기“제한”이 “구속”처럼 들릴 수 있지만, 사실\textbf{AI를 당신에게 가장 적합한 출력으로 수렴시키는}강력한 도구입니다. 일반적인 차원으로는 글자 수, 형식, 어조, 회피 항목, 필수 포함 항목이 있습니다.실용 예시:“\textbf{200자를 넘기지 말고, 항목별로 나열}.”(길이+형식)“\textbf{친근하고 자연스러운 어조로, 친구와 대화하듯이}, '첫째, 둘째, 셋째' 같은 관용구를 사용하지 마세요.”(어조)“\textbf{전문 용어를 사용하지 마세요}, 전문 용어를 만나면 함께 설명해주세요.”(용어 제한)“\textbf{각 단락은 3줄을 넘기지 마세요}, 스마트폰에서 읽기 편하도록.”(편집)“\textbf{반드시 포함: 사용 시나리오, 장단점, 내 예산 적합도}.”(필수 포함 항목)“\textbf{구체적인 투자 조언을 주지 마세요}, 원리만 설명해주세요.”(경계)\textbf{왜 효과적인가?}제한이 없으면, AI는 “자유롭게 발휘”하기 쉽습니다 — 너무 학술적이거나, 너무 길거나, 주제를 벗어남. 명확한 제한 = 명확한 “전달 기준”, 양측이 기대치를 일치시키고, 결과가 통제 가능합니다.5. 프롬프트 품질 평가: 좋은 프롬프트의 5가지 특징좋은 프롬프트를 쓰는 것은 일회성이 아니며,\textbf{다 쓴 후 먼저 이 5가지를 자체 점검}한 다음 AI에게 보내세요:\textbf{구체적이지 추상적이지 않게}: “글을 써줘”가 아니라, “800자, 3개 섹션, 초보 엄마를 대상으로 한 스마트폰 구매 가이드를 써줘”.\textbf{배경을 포함하고 고립시키지 않기}: AI에게 “나는 누구인지, 왜 이 질문을 하는지” 알려주어야, “당신만을 위한” 답변을 줄 수 있습니다.\textbf{명확한 출력 형식}: 표, 목록, 단락, 아니면 JSON? 미리 말하세요.\textbf{제한이 제대로 도달하게}: 글자 수, 스타일, 회피 항목, 셋 중 하나만 선택해도 뚜렷한 개선이 있습니다.\textbf{AI가 검증할 수 있게}: 질문 후 AI에게 요구사항을 다시 말하도록 요청; 말로 설명 못 하면 제대로 말한 것이 아닙니다.6. 자주 묻는 질문과 해법증상해결책답변이 너무 광범위하고 실용적이지 않음배경 추가(당신이 누구인지/어떤 시나리오인지)답변이 너무 학술적/너무 공식적어조 제한 추가 + “공식적인 말” 목록 거부답변이 주제를 벗어남/질문에 답하지 않음작업을 동사로 다시 말하고, “이것만 답해” 요구답변이 너무 길어서 다 읽을 수 없음글자 수 강제 제한(200자 이내)AI가 진지한 척 거짓말을 함“확실하지 않으면 '확실하지 않음'이라고 표시해” 추가 + 중요한 사실 자체 점검첫 번째 버전이 만족스럽지 않음다시 쓰지 말고, 원래 프롬프트에 “조금 더 ...” 추가7. 세 가지 완전한 시나리오 사례시나리오 1: 생활 어시스턴트(아침 식사 조합 조언)\textbf{프롬프트}:당신은 영양 조합에 능한 아침 블로거입니다(\textbf{역할}). 저는 직장인 엄마로, 매일 아침 15분만 아침을 준비할 시간이 있고, 집에 6살 아이와 60살 시어머니가 있습니다(\textbf{배경}).\textbf{7일 동안 겹치지 않는 아침 식사 조합을 나열}해주세요(\textbf{작업}). 각 아침 식사는: (1) 준비 시간 15분 이내 (2) 영양 균형(단백질+탄수화물+채소와 과일) (3) 온 가족이 즐길 수 있음 (4) 재료는 일반 슈퍼마켓에서 구매 가능(\textbf{제한}).\textbf{왜 효과적인가}: 역할이 전문성을 정하고, 배경이 세 유형의 공통점을 제공하며, 작업이 명확한 결과물을 제시하고, 4가지 제한이 실행 가능성을 보장.시나리오 2: 학습 보조(영어 6급 시험 준비)\textbf{프롬프트}:당신은\textbf{경험 풍부한 영어 6급 선생님}입니다(역할). 저는 일반 학부생으로,\textbf{4급 480점, 6급 목표 500점, 3개월 남음}(배경).\textbf{실행 가능한 12주 시험 준비 계획을 세워주세요}(작업). 요구 사항: (1) 매주 청해독작의 시간 배분 명확히 (2) 매월의 중점 돌파 방향 제시 (3) 적합한 복습 자료 3권 추천(중국어로 이유 설명) (4) 표 형식으로 출력(제한).\textbf{왜 효과적인가}: 역할이 전문 배경을 정하고, 배경이 점수와 시간 기준을 제공하며, 작업이 “단계별 계획”을 명확히 하고, 제한이 “표 형식”과 “실행 가능”을 요구.시나리오 3: 업무 효율(회의록 정리)\textbf{프롬프트}:당신은\textbf{숙련된 행정 비서}입니다(역할). 저는 방금 1시간 부서 회의를 마쳤고,\textbf{Q2 마케팅 홍보 방안을 논의}했습니다(배경).\textbf{다음 회의 기록을 구조화된 회의록으로 정리해주세요}(작업). 요구 사항: (1) “결정 사항 / 할 일 / 위험과 이의” 세 부분으로 구분 (2) 할 일에 담당자와 마감 시간 표시(언급되지 않은 경우 “확인 필요”로 작성) (3) 총 글자 수 600자 이내로 제한 (4) markdown 형식으로 출력(제한).\textbf{왜 효과적인가}: 역할이 행정 관점을 정하고, 배경이 회의 주제를 제공하며, 작업이 구조화된 출력을 명확히 하고, 4가지 제한이 결과를 바로 재사용할 수 있게 함.고급 플레이법: AI를 당신의 “개인 비서”로 만들기이러한 고급 활용법을 시도해 보고, AI의 진정한 힘을 경험해 보세요:\textbf{다회 대화 최적화}: AI의 답변이 만족스럽지 않을 때, “이건 너무 전문적이야, 초등학생도 이해할 수 있는 말로 다시 한 번 설명해줘”라고 말할 수 있습니다.\textbf{체인 오브 쏘트(CoT)}: AI에게 “단계별로 사고”하라고 하면, 답변 품질이 보통 더 높음: “이 수학 문제를 단계별로 분석해줘”.\textbf{다중 AI 협업}: A로 초안 작성, B로 다듬기, C로 번역 — 조합하여 사용하면 효과가 더 좋습니다.\textbf{시스템 프롬프트}: AI 설정에서 자신의 정체성, 선호도, 금기 등을 사전에 설정하여, 답변이 더 수요에 부합하도록 합니다.일곱, AI의 한계와 도전: 맑은 판단력 유지AI는 강력하지만, 결코 “만능”이 아닙니다. 한계를 명확히 인식해야 더 잘 사용할 수 있습니다:\textbf{“환각” 문제}: AI는 진지한 척 거짓말을 할 수 있습니다(전문 용어로 Hallucination). 존재하지 않는 역사적 사건, 잘못된 인용, 가짜 통계 데이터를 지어냅니다. 중요한 정보는 반드시 사람이 확인해야 합니다.\textbf{지식 마감}: 대부분의 AI 학습 데이터에는 마감일이 있습니다(예: 2024년 초), 이후 일어난 일은 알지 못합니다.\textbf{추론 능력 제한}: 복잡한 다단계 추론, 수학 증명, 논리적 함정에서는 AI가 여전히 쉽게 실수합니다.\textbf{프라이버시 위험}: 클라우드 AI 서비스에서는 대화 내용이 모델 학습에 사용될 수 있습니다. 민감한 정보는 온디바이스 AI를 우선 사용해야 합니다.\textbf{편견과 가치관}: AI는 학습 데이터의 편견을 담고 있을 수 있으며, 독립적 사고를 유지해야 합니다.AI 사용 팁AI가 생성한 콘텐츠는 정확하지 않을 수 있습니다(“환각” 문제), 중요한 정보는 반드시 확인하세요.개인 프라이버시 정보(주민등록번호, 비밀번호, 은행 정보 등)를 AI에 입력하지 마세요.프롬프트(Prompt)를 잘 활용하세요: 더 구체적이고 명확한 설명일수록, AI의 답변 품질이 높아집니다.AI는 도구이지 권위가 아니며, 독립적 사고와 판단력을 유지하세요.핵심 의사결정(의료, 법률, 재무)은 AI에 의존하지 말고 전문가와 상담하세요.도전해 보세요: 세 가지 AI 실전 문제\textbf{도전 1: 복잡한 명령}: 음성 비서에게: “내일 아침 7시에 우산 가져가는 것 알림, 그리고 오늘 오후 3시 일정을 4시로 변경.” 한 번에 다단계 작업을 완료할 수 있는지 보세요.\textbf{도전 2: AI 사진 편집}: 행인이 있는 관광 사진을 찍고, “마법 제거” 또는 “AI 지우개” 기능을 사용하여 행인을 깨끗이 제거. AI의 배경 “상상” 효과를 관찰.\textbf{도전 3: 프롬프트 최적화}: 먼저 간단한 프롬프트 “여행 가이드 한 단락 써줘”로 AI에게 물어보고, 다시 CRISPE 프레임워크로 자세한 프롬프트를 작성. 두 번의 답변 품질 차이를 비교.사례 2: 기기 변경 데이터 마이그레이션 완전한 흐름\textbf{시나리오}: iPhone 13 샤오미 15 브랜드 간 기기 변경.\textbf{단계}: (1) 구 기기 iCloud/위챗 자체 백업 (2) Apple ID 로그아웃 + “찾기” 끄기 (3) 새 기기에서 “샤오미 기기 변경” 앱으로 QR 코드 스캔 (4) 마이그레이션 항목 선택 (5) 위챗은 자체 “채팅 기록 마이그레이션”으로 별도 처리 (6) 핵심 데이터 확인(주소록/앨범/위챗) (7) 구 기기 암호화 삭제.\textbf{예상 결과}: 데이터 손실 0, 구 기기를 안전하게 재판매 가능.}
